\documentclass[10pt,a4paper]{article}

% ============================================================
% PACKAGES
% ============================================================
\usepackage[top=0.5in, bottom=0.45in, left=0.7in, right=0.7in]{geometry}
\usepackage{titlesec}
\usepackage{enumitem}
\usepackage{hyperref}
\usepackage{xcolor}
\usepackage[T1]{fontenc}
\usepackage{charter}
\usepackage{parskip}
\usepackage{needspace}

% ============================================================
% COLORS (match raadh.me)
% ============================================================
\definecolor{heading}{HTML}{0B4F5C}
\definecolor{link}{HTML}{0B4F5C}
\definecolor{body}{HTML}{1A1A1A}
\definecolor{muted}{HTML}{5A6470}

% ============================================================
% HYPERLINKS
% ============================================================
\hypersetup{
    colorlinks=true,
    linkcolor=link,
    urlcolor=link,
    pdfauthor={Rokib Al Dhin Raadh},
    pdftitle={Rokib Al Dhin Raadh - Curriculum Vitae},
    pdfsubject={AI researcher; founder and CEO of OXIEDO; inventor of ORMAS},
    pdfkeywords={Rokib Al Dhin Raadh, ORMAS, OXIEDO, DeepMath 2026, NeurIPS 2026, input-to-state stability, self-repairing neural networks},
}

% ============================================================
% SECTION STYLE
% ============================================================
\titleformat{\section}
  {\normalsize\bfseries\scshape\color{heading}}
  {}
  {0em}
  {}
  [\vspace{-1pt}{\color{heading}\titlerule[0.4pt]}]
\titlespacing{\section}{0pt}{9pt}{5pt}

\pagestyle{empty}

% ============================================================
% CV ENTRY: date in the left gutter, body indented
% ============================================================
% A list (not a minipage) so long entries can break across pages.
\newcommand{\cvitem}[3]{%
  \begin{list}{}{%
    \setlength{\leftmargin}{2.35cm}%
    \setlength{\labelwidth}{2.1cm}%
    \setlength{\labelsep}{0.25cm}%
    \setlength{\itemindent}{0pt}%
    \setlength{\listparindent}{0pt}%
    \setlength{\rightmargin}{0pt}%
    \setlength{\topsep}{0pt}%
    \setlength{\partopsep}{0pt}%
    \setlength{\parsep}{0pt}%
    \setlength{\itemsep}{0pt}%
  }%
  \item[{\makebox[2.1cm][l]{\footnotesize\color{muted}#1}}]%
  #2
  #3
  \end{list}%
  \vspace{1.5pt}
}

% Tight bullet list
\newlist{bullets}{itemize}{1}
\setlist[bullets]{leftmargin=1em, itemsep=0pt, parsep=0pt, topsep=1pt, label=\textbullet}

% Numbered publication list
\newlist{pubs}{enumerate}{1}
\setlist[pubs]{leftmargin=2.35cm, labelsep=0.25cm, itemsep=2pt, parsep=0pt, topsep=1pt, label={\footnotesize\color{muted}[\arabic*]}}

\begin{document}
\color{body}

% ============================================================
% HEADER
% ============================================================
\begin{center}
  {\LARGE\bfseries\color{heading} Rokib Al Dhin Raadh}\\[3pt]
  {\small AI researcher \enspace$\cdot$\enspace Founder \& CEO, \href{https://oxiedo.com}{OXIEDO} \enspace$\cdot$\enspace Inventor of ORMAS \enspace$\cdot$\enspace Dhaka, Bangladesh}\\[3pt]
  {\footnotesize
    \href{mailto:raadh@oxiedo.com}{raadh@oxiedo.com}
    \enspace$|$\enspace
    \href{https://www.raadh.me}{raadh.me}
    \enspace$|$\enspace
    \href{https://orcid.org/0009-0003-1178-5296}{ORCID 0009-0003-1178-5296}
    \enspace$|$\enspace
    \href{https://oxiedo.com}{oxiedo.com}
  }
\end{center}
\vspace{-2pt}

% ============================================================
% DISTINCTIONS (the skim layer)
% ============================================================
\section*{Distinctions}
\begin{bullets}
  \item \textbf{Sole-authored paper accepted after double-blind review} at DeepMath 2026, the Conference on the Mathematical Theory of Deep Neural Networks (The Ohio State University); one of 41 accepted papers.
  \item \textbf{Reviewer, NeurIPS 2026} Workshop on Trustworthy AI for Good; listed among its 490 reviewers.
  \item \textbf{Invited speaker, Cohere Labs} Open Science Community, November 2026.
  \item \textbf{Highest-ranked application} in its 2026 review cycle, Cosmos Institute.
  \item \textbf{Inventor of ORMAS}, a self-repairing neural network: recovers to 80.3\% accuracy after trained layers are zeroed mid-training, while standard training stays at chance (10.0\%).
  \item \textbf{Founder \& CEO of OXIEDO} (2023); accepted to 1752vc Ignite from the top 1\% of applicants.
\end{bullets}

% ============================================================
% RESEARCH
% ============================================================
\section*{Research}
\noindent{\small\textbf{Focus.} Stability theory of self-repairing training: neural networks that detect, repair and record their own failures while they train, analysed as switched control systems (input-to-state stability).}

\vspace{3pt}
\cvitem{2024--now}{%
  \textbf{Independent researcher} \quad {\small ORMAS: self-repairing neural networks}
}{%
  \begin{bullets}
    \item Proved that, near a minimum, bounded repairs keep training within a computable radius of it (practical input-to-state stability), and that bounded repairs alone cannot revive a dead channel in a BatchNorm and ReLU layer: self-repairing training is a switched system of bounded corrections and isolated resets.
    \item Recovery after trained convolutional layers are zeroed: \textbf{80.3\%} against \textbf{10.00\%} (chance) for standard training, 3/3 seeds. Sequential learning: \textbf{94.6\%} retention against 47.3\% for \mbox{ResNet-18}.
    \item \textbf{383 controlled experiments} across four architecture families (DAG, two CNNs, \mbox{ResNet-18}; CIFAR-10/100), each reproducible from seed, all on one RTX~3090. 16,316 lines of PyTorch.
    \item Released \textbf{ormas-monitor} (PyPI), a training monitor that certifies dead channels under the theorem's conditions; on 22 released torchvision models it finds near-zero channels in 19 of the 21 that use BatchNorm--ReLU.
  \end{bullets}
}

% ============================================================
% PUBLICATIONS
% ============================================================
\section*{Publications}
\noindent{\footnotesize\color{muted}Peer-reviewed}
\begin{pubs}
  \item \textbf{R.~A.~D. Raadh}. ``Self-Repair as a Bounded Disturbance: Input-to-State Stability of Neural Network Training Dynamics.'' \textit{DeepMath 2026: Conference on the Mathematical Theory of Deep Neural Networks}, The Ohio State University, 2026. Poster; double-blind review. \href{https://openreview.net/forum?id=kowltBZEIv}{OpenReview} $\cdot$ \href{https://www.raadh.me/deepmath}{poster}
\end{pubs}
\noindent{\footnotesize\color{muted}Preprints and technical reports}
\begin{pubs}[resume]
  \item \textbf{R.~A.~D. Raadh}. ``ORMAS: Neural Architectural Transparency Enables Autonomous Self-Correction.'' Preprint, Zenodo, 2026, \href{https://doi.org/10.5281/zenodo.21730363}{doi:10.5281/zenodo.21730363}. 36-page supplement; code reproduces all 383 experiments.
  \item \textbf{R.~A.~D. Raadh}. ``OXIMO/AX09: Autonomous LLM-Dependent Commerce. A Twelve-Month Controlled Ablation Study.'' Technical report, 2026.
\end{pubs}

% ============================================================
% TALKS
% ============================================================
\section*{Invited Talks and Presentations}
\cvitem{Nov 2026}{%
  \textbf{Invited talk}, Cohere Labs Open Science Community.\\
  ``When the Update Is Not a Gradient: Input-to-State Stability of Self-Repairing Neural Network Training.''
}{}
\cvitem{Oct 2026}{%
  \textbf{Poster}, DeepMath 2026, The Ohio State University, Columbus, OH (poster session, 29 October).\\
  ``Self-Repair as a Bounded Disturbance: Input-to-State Stability of Neural Network Training Dynamics.''
}{}

% ============================================================
% SERVICE
% ============================================================
\section*{Professional Service}
\cvitem{2026}{%
  \textbf{Reviewer}, Trustworthy AI for Good Workshop, \textbf{NeurIPS 2026} (Paris). Reviewed submissions on mechanistic interpretability, attribution, auditing and post-deployment monitoring. \href{https://trustworthy-ai-for-good.github.io/\#reviewers}{Reviewer list}
}{}

% ============================================================
% VENTURES
% ============================================================
\section*{Ventures}
\cvitem{2023--now}{%
  \textbf{Founder \& CEO, \href{https://oxiedo.com}{OXIEDO}} \quad {\small Auditable AI training for regulated industries}
}{%
  \begin{bullets}
    \item Licenses ORMAS on-premise: every training run produces a tamper-evident record of what the model learned, so banks, hospitals and other regulated institutions can train on data they already hold.
    \item Accepted to 1752vc Ignite (top 1\% of applicants); offered a place by Freshmango; finalist, The Bridge (Entrepreneur First, San Francisco); invited to Onstage W26. In investment discussions with 500 Global and 1752vc.
  \end{bullets}
}
\cvitem{2023--2026}{%
  \textbf{Sole architect, OXIMO} \quad {\small Multi-agent operating system}
}{%
  \begin{bullets}
    \item 40,933 lines of production Python and 2,069 passing tests, rebuilt from a 106,000-line monolith (72\% smaller). Creates the specialist roles a task needs, with no human in the loop.
  \end{bullets}
}
\cvitem{2025--2026}{%
  \textbf{Founder \& CEO, Black Bloxie LTD} \quad {\small England \& Wales, company no.\ 16711223}
}{%
  \begin{bullets}
    \item Twelve-month controlled field study of OXIMO on a live e-commerce company: removing the system cut output by 91\%; restoring it recovered and exceeded the original level.
  \end{bullets}
}
\cvitem{2020--2024}{%
  \textbf{Five earlier ventures}, ages 12 to 17; built and sold a stock-prediction system for about \$10,000 at 15.
}{}

% ============================================================
% HONOURS
% ============================================================
\section*{Honours and Selection}
\cvitem{2026}{\textbf{IARCO 2026}, International Academic Research Competition: finalist; ranked 24th of 145 in the Junior division (\href{https://iarco.org/2026-winners}{results}).}{}
\cvitem{2026}{\textbf{Cosmos Institute}: ranked the ORMAS application highest in its review cycle; invited to its technical track.}{}
\cvitem{2026}{\textbf{Cohere Labs}: member, Open Science Community; the lab cited ORMAS for ``remarkable initiative in ML safety and auditability''.}{}
\cvitem{2026}{\textbf{MIT Open Learning}: offered a 40\% scholarship on a further programme.}{}

% ============================================================
% SOFTWARE
% ============================================================
\section*{Software}
\cvitem{2026}{\textbf{\href{https://pypi.org/project/ormas-monitor/}{ormas-monitor}} (PyPI): per-layer alarms, certified dead channels, optional repair and a tamper-evident training record for PyTorch. \href{https://oxiedo.github.io/ormas-monitor/}{Docs}}{}
\cvitem{2026}{\textbf{ORMAS reference implementation}: reproduces all 383 experiments. \href{https://anonymous.4open.science/r/ormas-EB73}{Source} $\cdot$ \href{https://drive.google.com/file/d/1CDaMIpTZ_8Mkot9D-O7JU29mDopq_Bdl/view}{run archive}}{}

% ============================================================
% EDUCATION
% ============================================================
\section*{Education}
\cvitem{now}{\textbf{Higher Secondary Certificate (HSC)}, final year, Rajuk Uttara Model College, Dhaka.}{}
\cvitem{2025--2026}{\textbf{MIT Open Learning (MITx)}: seven programmes completed, in foundation models, AI in medicine, energy, transportation and AI entrepreneurship. \textbf{Deep Learning Specialization} (five courses), DeepLearning.AI.}{}

% ============================================================
% OTHER
% ============================================================
\section*{Other}
\cvitem{Nov 2024}{Guinness World Record attempt: 28 unedited livestreamed working sessions over two weeks.}{}
\cvitem{}{Languages: English (fluent), Bengali (native).}{}

\vfill
\begin{flushright}
  \footnotesize\color{muted}\textit{Last updated November 2026}
\end{flushright}

\end{document}
